% ---------------------------------------------------------------------------
% Section: Training through time: the chained expansion, its applications,
%          and why the late stage is not predictable from the initialisation
%
% Drop-in section for the paper. Conventions of Part III: Pi = I, c_p = 1,
% gradient flow, E = \hat Y - Y, K(theta) the (1/n-normalised) tangent kernel
% so that \dot E = -K E, J the Jacobian on the data, theta_0 the initialisation,
% lambda_1 = lambda_max(K), H_2 the second-derivative bound of Part I.
% Cross-references to earlier results (Proposition 20 = the analyticity radius
% of expansions about theta_0; Section sec:held-rough = the held rough error)
% are hard-coded; replace them with \ref{...} to your labels. Environments:
% theorem, proposition, corollary, definition, assumption, remark, proof.
% Every numerical statement below was obtained at n = 200 training points on
% the Lorenz-63 testbed and is marked [n=200]; see the last subsection.
% ---------------------------------------------------------------------------

\section{Training through time: the chained expansion and the limits of prediction}
\label{sec:chain}

Proposition~20 says that no series in $\theta-\theta_0$ reaches the stretch of
training in which the rough part of the error is removed, and
Section~\ref{sec:held-rough} replaced the expansion variable by the peak
feature roughness. This section keeps $\theta$ as the variable and moves the
expansion point instead: training is written as a chain of expansions, each
about the current state and each valid within its own radius, whose
composition is exact. The chain gives three kinds of statement. Identities
that hold at every time for any network in a broad class (\S\ref{ssec:chain-identities});
structure theorems for what gradient flow does to the weights and the kernel
at each stage, which is where the low-rank ``spikes'' and the targeted growth
of the kernel come from (\S\ref{ssec:chain-spikes}--\ref{ssec:chain-richness});
and applications in which an open question becomes a computation with the
chain or a theorem about its leading orders (\S\ref{ssec:chain-apps}). The
last two subsections state why the late stage cannot be predicted from the
initialisation without carrying the state forward, in a form that is a
theorem rather than an impression (\S\ref{ssec:chain-impossible}), and what a
statistical description can and cannot promise (\S\ref{ssec:chain-statistical}).

\subsection{Setting and exact identities}
\label{ssec:chain-identities}

\paragraph{Model class.} $N_\theta$ is any composition of \emph{parameter-linear
blocks} $z_\ell = W_\ell a_\ell$, where $a_\ell$ is whatever the block receives
(dense, convolution, attention projections, normalisation gains), and
\emph{parameter-free maps} between them (pointwise activations, softmax,
normalisation), differentiable almost everywhere in $\theta$. ReLU qualifies.
Training is gradient flow on the squared loss, or any update of the
data-indexed outer-product form below (GD, SGD on a subset of the data,
momentum through its integral); Adam is treated in \S\ref{ssec:chain-apps}.
Write $\Psi(t)=J(\theta(t))\in\mathbb R^{N\times P}$ for the Jacobian on the $n$
data points ($N=n\,n_y$), $K=\Psi\Psi^{\!\top}/n$, $\delta_\ell^c(x)=\partial
\hat Y_c(x)/\partial z_\ell$ for the backprop signal into block $\ell$, and
$U(t)=\int_0^t E$ for the integrated residual.

\begin{proposition}[Four identities]
\label{prop:chain-identities}
For every $t$ at which $N_\theta$ is differentiable:
\begin{enumerate}
\item[(a)] $\dot E=-K(t)E$; hence $\tfrac{d}{dt}\|E\|^2=-2\langle E,KE\rangle$
and $E(T)=\mathrm{T}\!\exp\!\big(-\!\int_0^T K\big)E(0)$.
\item[(b)] $K=\frac1n\sum_\ell A_\ell\odot D_\ell$, with $A_\ell[i,j]=a_\ell(x_i)\!\cdot\!a_\ell(x_j)+1$
the Gram of the block's inputs and
$D_\ell[(i,c),(j,d)]=\sum_q\delta^c_{\ell,q}(x_i)\,\delta^d_{\ell,q}(x_j)$ the Gram
of its backprop signals.
\item[(c)] $W_\ell(t)-W_\ell(0)=-\frac1n\int_0^t\sum_i\sum_c\delta_\ell^c(x_i)\,E_c(x_i)\,a_\ell(x_i)^{\!\top}$:
the change of every block is a sum of outer products indexed by the data, its
row space lies in the span of the block's inputs on the data, and every block is
driven by the same residual.
\item[(d)] For a unit $q$ of block $\ell$ with preactivation $z_q$ and incoming
backward signal $m_q=W_{\ell+1}^{\!\top}\delta_{\ell+1}$,
\[
\dot z_q=\dot W_\ell a_\ell+W_\ell\,\dot a_\ell,\qquad
\dot m_q=\dot W_{\ell+1}^{\!\top}\delta_{\ell+1}+W_{\ell+1}^{\!\top}\dot\delta_{\ell+1}.
\]
The first terms are given by (c) (the \emph{coherent} part); the second carry
the motion of the neighbouring blocks through the current weights (the
\emph{propagation} part). Together with (a) and (b) this is a closed system on
the data.
\end{enumerate}
\end{proposition}
\begin{proof}
(a) is the chain rule for $\theta'=-\frac1nJ^{\!\top}E$. (b): $\partial\hat
Y_c(x)/\partial W_\ell[q,p]=\delta^c_{\ell,q}(x)\,a_{\ell,p}(x)$. (c): integrate
the gradient. (d): differentiate the definitions.
\end{proof}

\begin{remark}
Everything in this section is a consequence of
Proposition~\ref{prop:chain-identities}. In particular the propagation part of
(d) is the only place where the initialisation enters other than through the
unit statistics; omitting it (as any model built from the kernel alone must)
was measured to cost $11$--$28\%$ of the unit displacements early and a factor
$2$ in the late error on the reference problem [n=200].
\end{remark}

\subsection{The stage identity and the chain}
\label{ssec:chain-stage}

Fix stage boundaries $0=t_0<t_1<\dots<t_S=T$ and write
$\delta\Psi_s=\Psi(t_s)-\Psi(t_{s-1})$.

\begin{proposition}[Stage identity]
\label{prop:stage}
$K(t_s)-K(t_{s-1})=\mathrm{Cross}_s+Q_s$ with
$Q_s=\delta\Psi_s\delta\Psi_s^{\!\top}/n\succeq0$ and
$\mathrm{Cross}_s=(\Psi\delta\Psi_s^{\!\top}+\delta\Psi_s\Psi^{\!\top})/n$,
$\Psi=\Psi(t_{s-1})$. For a unit vector $v$, the quadratic form changes by
$\frac1n\big[2\langle\Psi^{\!\top}v,\delta\Psi_s^{\!\top}v\rangle+\|\delta\Psi_s^{\!\top}v\|^2\big]$,
which is positive whenever $\|\delta\Psi_s^{\!\top}v\|>2\|\Psi^{\!\top}v\|=2\sqrt{n\lambda_v}$.
\end{proposition}

\begin{theorem}[Feature-sum identity and bracket at $T$]
\label{thm:bracket}
$\Psi_T=\Psi_0+\sum_s\delta\Psi_s$, so for every unit $v$, with
$D_v=\sum_s\delta\Psi_s^{\!\top}v$ and $\ell_v=\|\Psi_0^{\!\top}v\|$,
\[
\tfrac1n\big(\|D_v\|-\ell_v\big)_+^2\;\le\;v^{\!\top}K_Tv\;\le\;\tfrac1n\big(\|D_v\|+\ell_v\big)^2 .
\]
$D_v$ splits exactly into coherent and propagation parts by
Proposition~\ref{prop:chain-identities}(d). On the directions the residual has
driven, $\ell_v$ is negligible and the lower bound is the value: measured
$0.90$--$1.00$ of $v^{\!\top}K_Tv$ for every rank $\ge60$ on the reference
problem [n=200]. Where $J$ is $C^1$ along the path,
$|\lambda_k(K_T)-\lambda_k(K_0)|\le\int_0^T C(s)\|E(s)\|\,ds$ with
$C=\frac2{n^2}\|\Psi\|_2^2\|\Psi'\|_2$ (Weyl), so every lifted eigenvalue at $T$
is bracketed between two functionals of the residual history.
\end{theorem}

\begin{theorem}[Stage error and the chain]
\label{thm:chain}
Where $J$ is $C^2$ along the path,
$\delta\Psi_s=\alpha_s\dot\Psi(t_{s-1})+R_s$ with
$\|R_s\|\le\frac{\alpha_s^2}2\sup\|\ddot\Psi\|$; choosing
$\alpha_s\|\ddot\Psi\|\le\varepsilon\|\dot\Psi\|$ makes each stage's first-order
increment accurate to relative $\varepsilon$, and the composition of the
stages is the exact trajectory. The quantity $\|\ddot\Psi\|/\|\dot\Psi\|$ is the
inverse of the stage length over which the linearisation at the current state
is trustworthy; its norm along the descent direction is $H_2$.
\end{theorem}

\subsection{What a stage does to the weights: the spike theorem}
\label{ssec:chain-spikes}

For a hidden block whose outgoing weights are $W_{\ell+1}$, write
$S[i,q]=\sigma'(z_q(x_i))$ for the local derivative of the map following the
block, $M[i,q]$ for the backprop coefficient contracted with the drive (for the
last hidden block $M=U\,W_{\mathrm{out}}$, of rank $\le n_y$), and $A$ for the
block's inputs on the data. By Proposition~\ref{prop:chain-identities}(c),
$\Delta W=-\frac1n(S\odot M)^{\!\top}A$.

\begin{theorem}[Core and bulk]
\label{thm:spike}
\begin{enumerate}
\item[(a)] (\emph{at initialisation}) Split $S=h_0+(S-h_0)$ with $h_0$ the mean of
$\sigma'$ over the data. Then $\Delta W=\Delta W^{\mathrm{core}}+\Delta W^{\mathrm{bulk}}$ with
$\Delta W^{\mathrm{core}}=-\frac{h_0}nM^{\!\top}A$ of rank $\le n_y$, its left
singular space spanned by the rows of $W_{\mathrm{out}}$ and its right singular
space by $A^{\!\top}U$. If the block's input weights are independent of
$W_{\ell+1}$ (true at $\theta_0$), the expected bulk-to-core energy ratio is
$\rho=\mathrm{Var}(\sigma'(z))/\mathbb E[\sigma'(z)]^2$ at the preactivation scale.
The output block has rank $\le n_y$ and the input block rank $\le d_{\mathrm{in}}+1$ exactly.
\item[(b)] (\emph{any time, no independence}) For any rank-$r$ approximation
$S_r$ of $S$ with remainder $\tilde S$,
$\mathrm{rank}\big((S_r\odot M)^{\!\top}A\big)\le n_y r$ and
$\|(\tilde S\odot M)^{\!\top}A\|_F\le\|\tilde S\|_{\max}\|M\|_F\|A\|_2/n$. Since
$Z=A\,W_\ell^{\!\top}+b$, the numerical rank of $S=\sigma'(Z)$ is controlled by the
spectrum of the block's input Gram and by the Hermite coefficients of
$\sigma'$ at the current preactivation scale, both quantities the chain
carries; for deeper blocks the bound composes multiplicatively down the stack.
\end{enumerate}
\end{theorem}
\begin{proof}
(a) Expand $\sigma'$; the core is $h_0M$; the bulk's expected energy factorises
by independence of $S_{iq}$ and $M_{iq}$. (b) $\mathrm{rank}(X\odot Y)\le
\mathrm{rank}X\cdot\mathrm{rank}Y$, the entrywise Frobenius bound, the
column-space inclusion for $Z$, and the Hermite expansion.
\end{proof}


\subsection{The richness expansion, its radius, and the chain as its resummation}
\label{ssec:chain-richness}

Let $\epsilon$ be any parameter that multiplies the weight velocity at fixed
residual dynamics in a suitable time: the inverse output scale $1/s$ with
$\tau=s^2t$ for centred outputs $\hat Y=s(N_\theta-N_{\theta_0})$; the Euler step
$\Delta$ of \S\ref{ssec:chain-apps}; a $\mu$P-type multiplier; an inverse-width
factor. Then
\[
\frac{dE}{d\tau}=-K(\theta)E,\qquad\frac{d\theta}{d\tau}=-\frac{\epsilon}{n}J(\theta)^{\!\top}E .
\]

\begin{theorem}[All-orders richness expansion]
\label{thm:richness}
For analytic parameter-free maps and each fixed $\tau$ there is $R(\tau)>0$ such
that every analytic observable of the state is a convergent power series in
$\epsilon$ for $|\epsilon|<R(\tau)$; in particular
\[
\frac{v^{\!\top}K(\tau)v}{v^{\!\top}K(0)v}-1=\sum_{k\ge1}a_k(\tau)\,\epsilon^k ,
\]
with coefficients determined recursively from $\theta_0$ and the lazy
trajectory alone: $a_0$ is kernel regression with $K(\theta_0)$,
$\theta^{(1)}(\tau)=-\frac1nJ_0^{\!\top}U_{\mathrm{lazy}}(\tau)$,
$a_1=2\langle\Psi_0^{\!\top}v,J'[\theta^{(1)}]^{\!\top}v\rangle/(n\lambda_v)$, and $a_k$
involves derivatives of $J$ at $\theta_0$ up to order $k$ along lower-order
displacements and the lower-order residual corrections (linear inhomogeneous
equations, variation of constants in the eigenbasis of $K_0$). No quantity along
a trained path enters at any order.
\end{theorem}

\begin{theorem}[Radius, and the chain as resummation]
\label{thm:radius}
$R(\tau)>0$ by analyticity. Within the region in which the stage criterion of
Theorem~\ref{thm:chain} holds at $\theta_0$ the first-order term of the series
is accurate to the stated relative error, and along a direction $v$ the
first-order lift is a valid approximation while
$\|\delta\Psi^{\!\top}v\|<\|\Psi_0^{\!\top}v\|$ (Proposition~\ref{prop:stage}); these
are truncation statements and do not bound $R(\tau)$ from below, which would
need a majorant estimate on the derivatives of $J$ at $\theta_0$ (available
for analytic activations on a ball; not carried out). Where the lift along
$v$ reaches order one the truncated series fails sharply in the measurements,
consistent with the singularity of $\epsilon\mapsto$ trajectory sitting there.
Re-expanding at each
stage boundary (Theorem~\ref{thm:chain}) continues the series past $R(\tau)$
by moving the expansion point along the path; within $R(\tau)$ the two agree,
beyond it only the chain is defined.
\end{theorem}
\begin{proof}
Analytic dependence of ODE solutions on parameters; the recursion is the Taylor
expansion of the vector field in $\epsilon$; the truncation statements are
Proposition~20 at the expansion point $\theta_0$ together with
Proposition~\ref{prop:stage}; the resummation uses only $R(\tau)>0$.
\end{proof}


\subsection{Applications}
\label{ssec:chain-apps}

\paragraph{Euler versus direct.} Let $y_D=\Phi_\Delta(x)$ and
$y_E=(\Phi_\Delta(x)-x)/\Delta$ be the direct and the Euler target, trained from
the same $\theta_0$ with centred outputs, and write the direct residual on the
increment as $E_D=\Delta\tilde E$.
\begin{theorem}
\label{thm:euler}
$\tilde E$ and $E_E$ obey the same residual equation with their own kernels,
start from the same value, and the weight velocities satisfy
$\dot\theta_D=-\frac\Delta nJ^{\!\top}\tilde E$ against
$\dot\theta_E=-\frac1nJ^{\!\top}E_E$: the direct problem is the Euler problem with
the richness parameter multiplied by $\Delta$ (Theorem~\ref{thm:richness}
evaluated at a smaller argument). Hence (i) direct's stage lifts on the
increment are $\Delta^2$ times Euler's at equal normalised residual, and (ii) as
$\Delta\to0$ at fixed budget the direct network is in the lazy regime on the
increment, after a $\Delta$-independent early phase in which it fits the
identity component of its target, and its normalised error converges to the
frozen-kernel error of the increment target.
\end{theorem}
For direct as $\Delta\to0$ the held-rough statements of
Section~\ref{sec:held-rough} are therefore frozen-kernel statements (no lift,
fixed basis); they are not tight, since the floor decays at the rate set by
the top of the rough subspace while the frozen-kernel error decays mode by
mode. Euler's error is the full feature-learning error, below the frozen-kernel
value by the feature-learning gain, which grows with the budget and with the
increment's roughness; this comparison is conditional on the rate comparison
below. The static argument of the previous sections (``the increment is a smoother
target'') is the frozen-kernel shadow of this: dividing by $\Delta$ multiplies the
drive by $1/\Delta$ and the lifts by $1/\Delta^2$, so Euler trains rich on the
increment while direct trains lazy on it.

\paragraph{Progressive sharpening, the initial reduction phase, and the edge of stability.}
By Proposition~\ref{prop:stage},
$\dot\lambda_1=v_1^{\!\top}(\mathrm{Cross}+Q)v_1$ per stage with $Q\succeq0$. Along
the residual the cross term at $t=0$ equals $-\frac{2\alpha}{n^2}\langle
E,N''_{\theta}[g,g]\rangle$ with $g=J^{\!\top}E$, exactly, and it is negative on
$90$--$91\%$ of the eigen-directions for the raw residual $E=\hat Y_0-Y$ and on
$48$--$55\%$ (no systematic sign) for the centred residual $-Y$ [n=200]: the
initial sharpness-\emph{reduction} phase reported in the warmup literature is
the network unlearning its own initial function and is absent for centred
initialisations. Once learning has accumulated, the coherent part of the cross
term is the inner product of the accumulated core above with the stage core
below through the shared drive, positive while the drive keeps its direction
(measured positive on $80$--$100\%$ of modes from $t\approx1$ on): progressive
sharpening, faster with depth by composition ($5\times,11\times,20\times$
top gain at depths $2,3,6$ on the $\Delta=0.5$ target and $86\times$ at depth $6$
on $\Delta=1.0$ [n=200]). Taking GD steps as stages, above the
stability threshold the top component of the residual alternates in sign, the
stage cores along it alternate, their cross terms with the accumulated core
are negative at the core level, and $\lambda_1$ decreases until the step is
stable again: sharpening and self-stabilisation are the two signs of one term.
What is not proved is the sign of the init-feature part of the cross term
after the transient (Lemma~\ref{lem:sign} below).

\paragraph{Adam.} For $\dot\theta=-P(t)\nabla L$ with $P$ diagonal positive, (a)
of Proposition~\ref{prop:chain-identities} holds with $K_P=\Psi P\Psi^{\!\top}/n$,
(c) with the outer products multiplied entrywise by $P_\ell$, and
Theorem~\ref{thm:bracket} at a fixed time with $P_T^{1/2}$-weighted features.
Along the path the stage identity acquires a third term: with
$\delta P_s=P(t_s)-P(t_{s-1})$,
$K_P(t_s)-K_P(t_{s-1})=\mathrm{Cross}_{P,s}+Q_{P,s}+\Psi(t_s)\,\delta P_s\,\Psi(t_s)^{\!\top}/n$,
and $\delta P_s$ is negative on every coordinate whose second-moment estimate
is still growing. The lift and the preconditioner's decay therefore compete,
the monotone-spectrum inputs below do not transfer to $K_P$ without a bound
on this term, and with momentum the drive is the averaged gradient, so the
energy identity has no sign. What transfers is the set of identities, the
fixed-time bracket and the spike statement; the discrete regime in which Adam
runs in practice (its own edge of stability) is outside gradient flow
altogether. The spike statement needs one more fact: $P$ is
built from the gradient block itself, so $P\odot G$ is a normalised sign
pattern of a rank-$\le n_y$ matrix, and the same mean-plus-fluctuation
argument as Theorem~\ref{thm:spike} gives its core fraction as the squared first
Hermite coefficient of the normaliser under the entry distribution of $G$:
$2/\pi$ for Gaussian-like entries, $1$ for coordinate-aligned (sparse) singular
vectors, with the time-averaged second moment lowering it further (measured
$0.68$--$0.81$, $0.95$--$1.00$, $0.52$ respectively on synthetic blocks). This
is the rotation-dependent property that distinguishes Adam from SGD.

\paragraph{Validation, and minibatches.} With $J_{\mathrm{val}}$ the Jacobian on validation
points, $\dot{\hat Y}_{\mathrm{val}}=-K_{vt}E$ with
$K_{vt}=J_{\mathrm{val}}J^{\!\top}/n$, exactly, so
$-\frac d{dt}\frac12\|E_{\mathrm{val}}\|^2=E_{\mathrm{val}}^{\!\top}K_{vt}(0)E+
E_{\mathrm{val}}^{\!\top}[K_{vt}(t)-K_{vt}(0)]E$: a lazy-transfer part and a
feature-learning-transfer part. The coherent part of the lift on validation
points is the training stage core transported by the cross activation Gram;
under equal unit statistics on the two samples the whole chain applies with
that substitution. The same identity covers minibatch training: a step on a
batch $B$ drives the residual on the full set by $\dot E=-K_{\mathrm{full},B}\,g_B$,
which is the validation identity with $B$ as the training set, so the chain
is exact for SGD with the drive history $\{g_{B_s}\}$, while the full-set
energy, like the validation energy, has no sign step by step.
Measured [n=200]: early both parts positive and the
validation rate $70$--$110\%$ of the training rate; late the lazy part
negative (the frozen kernel maps a residual living in sample-specific rough
modes into directions that increase validation error, the static held-rough
mechanism seen as a sign) while the lifted part stays positive. At the sample
sizes of interest, where training and validation agree, the sign pattern is a
prediction to be checked rather than a result.

\paragraph{Delayed feature learning (``grokking'').} The mechanism is
Theorems~\ref{thm:spike}--\ref{thm:bracket}: the residual drives a positive
lift along the directions it has visited, so a direction too slow for the
initial kernel to fit within the budget is fit later through the lift. The
onset is Theorem~\ref{thm:richness} with $\epsilon=1/s$: an init-computable
series, verified to $0$--$16\%$ in its range, with a proved radius; the delay
to a small lift is its inversion, and because $a_1$ changes sign no universal
exponent in $s$ exists, which accounts for the disagreement among reported
scalings. The event itself, a lift of order one or more along a misaligned
direction, lies outside $R(\tau)$ by construction and is a chain computation;
by \S\ref{ssec:chain-impossible} it is not predictable from the
initialisation in any other way. Rigorous timings in the literature (weight
decay, modular arithmetic, two-layer mean field) are special cases evaluated
inside their own radii.

\paragraph{The error at $T$, term by term.} Write $K(t)=K_0+\Delta(t)$ with
$\Delta$ the accumulated stage increments. The residual equation has the exact
Dyson form
\[
E(T)=\sum_{m\ge0}(-1)^m\!\!\int_{T>u_1>\dots>u_m>0}\!\! e^{-K_0(T-u_1)}\Delta(u_1)\,e^{-K_0(u_1-u_2)}\cdots\Delta(u_m)\,e^{-K_0u_m}E_0,
\]
absolutely convergent for bounded $\Delta$ (the $m$-th term is bounded by
$(\int\|\Delta\|)^m/m!\,\|E_0\|$), so the series is the solution, not an
approximation. Every term is a product of lazy propagators and increments, and
every increment is built from the stage cores of Theorem~\ref{thm:spike}, so
$E(T)$ lies in the span generated by the lazy flow of $E_0$ and the stage
targets. The first-order term paired with the lazy solution is
$-\int_0^T\langle E_{\mathrm{lazy}}(2T-u),\Delta(u)E_{\mathrm{lazy}}(u)\rangle\,du$,
whose sign is fixed by two properties at each $u$: $\Delta(u)\succeq0$ as an
operator on the subspace spanned by the lazy residual at $u$ and at $2T-u$
(the sign lemma below) and the alignment of
$E_{\mathrm{lazy}}(u)$ with $E_{\mathrm{lazy}}(2T-u)$ in the modes $\Delta$ touches
(persistence; automatic in a fixed basis). Under both,
$\|E(T)\|^2\le\|E_{\mathrm{lazy}}(T)\|^2-2\int_0^T\langle
E_{\mathrm{lazy}}(2T-u),\Delta(u)E_{\mathrm{lazy}}(u)\rangle du+R_2$, with an
explicit nonnegative first correction and a remainder bounded by
$\|E_0\|^2(D^2+2e^Dr+r^2)$, $D=\int_0^T\|\Delta\|$, $r=e^D-1-D$, from the Dyson
bound on the terms $m\ge2$. The remainder is not sign-controlled by the two
properties once the basis rotates; in a fixed basis the series resums
exactly to the mode-wise bound of Section~\ref{sec:held-rough}, and there the
two properties are the only hypotheses to all orders. With rotation the
first-order inequality is what the expansion proves, and the all-orders
comparison is organised differently below.

\paragraph{What the data say about the two hypotheses.} On the reference run
every initial eigen-direction carrying the late residual satisfies the
feature-change inequality of Theorem~\ref{thm:bracket}, so the accumulated
increment is positive along each of them [n=200]. This is evidence for the
sign condition on that subspace, not a proof of it: positivity of a quadratic
form on each of a set of directions does not give positivity on their span,
the condition the theorems use is the smallest eigenvalue of the increment
compressed to the subspace, and the monotone-spectrum statements below need
it for each stage increment, not only for the accumulated one. What fails
outright is persistence in the initial
basis: the subspace carrying the residual at $t=44$ is only $55\%$ retained at
$t=428$, the accumulated increment on those modes being $10^2$--$10^3$ times the
initial gaps. Persistence is therefore false there, not merely unproved, and
the error statement must be organised in the instantaneous eigenframe, where
the residual's rank profile moves slowly (the front of
Section~\ref{sec:held-rough}) even though its direction does not; the
fixed-basis bounds are then the adiabatic limit and the rotation enters as a
named remainder, coupling over current gap, per stage. That derivation is the
remaining step.

\paragraph{The moving-frame identity.} With $K(t)=V\Lambda V^{\!\top}$ and
$c=V^{\!\top}E$ the residual in the instantaneous eigenframe,
$\dot c=-\Lambda c-Ac$ with $A=V^{\!\top}\dot V$ antisymmetric, so
$\frac d{dt}\|c\|^2=-2c^{\!\top}\Lambda c$ exactly: rotation redistributes energy
among the current modes and does not change the norm. If every stage
increment from $t_0$ on is positive semidefinite on a subspace containing the
residual, then $\langle E,K(t)E\rangle\ge\langle E,K(t_0)E\rangle$ and the descent
rate is at least the frozen-kernel rate on the \emph{current} residual; if the
increments are positive on the whole space, the sorted eigenvalues are
nondecreasing (Weyl) and the same reads, in the moving frame, as the frozen
sorted spectrum evaluated on the current rank profile. Neither form needs
eigenvector persistence; both need per-stage positivity, which is the sign
lemma below stated per stage and is a hypothesis, checked on the data by the
smallest eigenvalue of the compressed stage increment.
What a closed comparison with the lazy trajectory then needs is the evolution
of the rank profile, i.e.\ the moving cutoff of
Section~\ref{sec:held-rough}: in the moving frame persistence is a statement
about where the residual sits, not which way it points, and the fixed-basis
bounds are the case $A=0$.

\paragraph{The cutoff bound.} Fix the index set $u=\{k:\lambda_k(T)T<1\}$ of
modes still unfit at $T$, with cutoff rank $k^*=\min u^c$, and transport it
along the eigenvector flow by continuity (through crossings, which are
massive here: sorted rank is not an invariant label), $P(t)=V(t)\Pi_uV(t)^{\!\top}$. With $E_u=PE$, $E_f=(I-P)E$
and the cross-front rotation rate $\rho=\|A_{uf}\|_2$, the spectral norm of
the $u\times f$ block of $A$,
\[
\tfrac d{dt}\|E_u\|^2=-2E_u^{\!\top}KE_u+L,\qquad \tfrac d{dt}\|E_f\|^2=-2E_f^{\!\top}KE_f-L,\qquad |L|\le2\rho\|E_u\|\|E_f\|,
\]
the leakage $L=2c_u^{\!\top}A_{uf}c_f$ being the only coupling between the two subspaces. Hence, with
$\lambda_u,\lambda_f$ the extreme eigenvalues of $K$ on the two tracked subspaces at each time,
\[
\|\Pi_uc(t_0)\|\,e^{-\int_{t_0}^T\lambda_u}-\!\int_{t_0}^T\!\rho\|E_f\|\;\le\;\|E(T)\|,\qquad
\|E(T)\|^2\le\|\Pi_uc(t_0)\|^2+2\!\int_{t_0}^T\!\rho\|E_u\|\|E_f\|+\|E_f(T)\|^2,
\]
with $\|E_f(T)\|\le\|(I-\Pi_u)c(t_0)\|e^{-\int\lambda_f}+\int\rho\|E_u\|$. At $t_0=0$,
$\Pi_uc(0)$ is the energy of the initial error on the init-time eigenvectors
that evolve into the unfit set at $T$ (not the energy above sorted rank $k^*$
at $t=0$, which mislabels the modes by an order of magnitude here), the cutoff
being set by the \emph{lifted} spectrum at $T$; computing it requires the
continuously tracked eigenvectors, which the chain carries. Three things
should be said about the inputs. A continuously tracked eigenvalue is not
monotone even when the sorted spectrum is, so $\int\lambda_u$ is a run
quantity; if the tracked eigenvalues in $u$ are nondecreasing, then
$\int_{t_0}^T\lambda_u<1$ and the floor is at least $e^{-1}\|\Pi_uc(t_0)\|$ minus
the leakage, which is the bracket's real content. For a generic path the
eigenvalues do not cross but avoid each other, and at an avoided crossing
between a mode in $u$ and a mode in $f$ the coupling over gap is large; those
events are exactly where $\rho$ is large, so the order-of-magnitude
mislabelling seen with sorted ranks does not vanish with continuity
transport but moves into the leakage integral, which has not been evaluated at
any sample size. And rotation among unfit modes, where the term-wise
corollary (iii) below places most of it, moves energy within $P$ and costs
nothing, while rotation across the front is weighted by the fit residual,
which is small except at the crossing events themselves. No persistence
hypothesis appears; the inputs are the tracked cutoff set, the leakage, and
per-stage positivity of the increments on the residual's subspace for the
decay rates. An equivalent formulation defines the unfit set at each time by
the current spectrum, $u(t)=\{k:\lambda_k(t)t<1\}$, the moving cutoff of
Section~\ref{sec:held-rough}; then $P$ jumps when a mode's $\lambda_kt$ crosses
one and the leakage carries those jumps explicitly. The cutoff rank is bracketed two-sidedly by
Theorem~\ref{thm:bracket} through the accumulated feature change, $\rho$ is the
per-stage coupling over current gap across $k^*$, and the static held-rough
floor of Section~\ref{sec:held-rough} is the case in which the cutoff is read
from the frozen spectrum: the lift moves the cutoff up in rank, and that is the
whole effect of feature learning on the error, in one scalar. For a direction
$v$, the time it leaves $u$ solves $\lambda_v(T_v)T_v=1$ and is bracketed by
$(\|D_v(t)\|\mp\ell_v)^2t/n=1$: the delayed-feature-learning event as an exact
crossing condition with a two-sided bracket from the accumulated drive and a
sharpness set by $\rho$ at the crossing.

\paragraph{The rate comparison, and the one lemma it rests on.}
\begin{lemma}[Sign of the cross term; status: derived for the core part, measured for the rest]
\label{lem:sign}
$v^{\!\top}\mathrm{Cross}_sv=\frac2n\big[\langle\Psi_0^{\!\top}v,\delta\Psi_s^{\!\top}v\rangle+\sum_{r<s}\langle\delta\Psi_r^{\!\top}v,\delta\Psi_s^{\!\top}v\rangle\big]$.
The second term's core-core part is nonnegative while the integrated residual
keeps its direction (shared drive); the first term is measured negative at
$t=0$ for raw residuals and sign-free for centred ones, positive after the
transient, and is not derived.
\end{lemma}
If the lemma holds from $t_0$ on, then $\dot L(t)\le-\langle E,K(t_0)E\rangle$ for
$t\ge t_0$: feature learning never slows the descent below what the kernel at
$t_0$ would produce on the current residual. A closed-form bound on
$\|E(T)\|$ needs in addition a persistence condition on the residual
direction, because $K(t)\succeq K(t_0)$ does not control mode-wise decay when
the eigenvectors rotate, and they rotate by about half in the relevant band by
mid-training [n=200].

As a theoretical object the chain is exact: infinitely many expansion points, each
to all orders, is the flow itself, and everything above uses only that. How a
truncation at one stage reaches $T$ (the error accounting that makes every
reduced version of the chain a controlled approximation), and the numerical
verifications of the identities and of the leading orders, are collected in
Appendix~\ref{app:chain}.

\paragraph{Term-wise corollaries.} Once the right series is used, several
statements become unconditional. (i) Direct is the richness series at argument
$\Delta\epsilon$, so its $k$-th lift term is $\Delta^k$ times Euler's: the lift
comparison holds order by order inside the radius with no sign lemma.
(ii) By the variational principle $\lambda_{\max}(K_0+\Delta)\ge\lambda_1+v_1^{\!\top}\Delta v_1$
and $\lambda_{\min}(K_0+\Delta)\le\lambda_{\min}+v_{\min}^{\!\top}\Delta v_{\min}$ for any
symmetric $\Delta$, so the totality of the higher-order terms pushes the top
eigenvalue up and the bottom down, exactly and with no convergence question
(the Rayleigh--Schr\"odinger series itself converges for $\|\Delta\|$ below half
the eigenvalue's gap, which the chain's stages satisfy for the top and only at
fine resolution for the tail); its second-order term
$\sum_{k>1}|v_1^{\!\top}\Delta v_k|^2/(\lambda_1-\lambda_k)\ge0$ is the leading piece: progressive sharpening of the
top mode has an unconditional second-order component, and the deep tail an
unconditional second-order depression, which is the measured early decrease of
the kernel along init-tail directions before the first-order lift overtakes it.
(iii) Mode $k$ rotates by $\sum_{j\ne k}(v_j^{\!\top}\Delta v_k)/(\lambda_k-\lambda_j)\,v_j$,
coupling over gap; with an init spectrum $\sim k^{-a}$ the relative gap is
$\sim a/k$, so the rotation grows with rank, i.e.\ with roughness: the
fixed-basis assumption fails exactly where the roughness theory places the
action, and $|\text{coupling}|/|\text{gap}|$ is a computable criterion for when a
direction keeps its identity. (iv) Two initialisations with the same unit
statistics differ in every Dyson term only through the propagation part of the
stage features, so the spread of $E(T)$ over initialisations is bounded by the
propagator times the accumulated propagation mismatch: the irreducible
uncertainty of any statistical description is an explicit, measurable,
width-dependent quantity, and whether the statistical chain becomes exact at
large width is the question whether that quantity shrinks with width.

\subsection{Consequences for the held rough error}
\label{ssec:chain-heldrough}

The results above replace the closure of Section~\ref{sec:held-rough} and
sharpen its exact part. (i) Proposition~\ref{prop:rough-energy} is the energy
identity projected on a fixed $P$; the cutoff bound is the same identity with
$P$ transported along the eigenvector flow, the static off-diagonal leak
becoming the cross-front rotation, and the fixed-$P$ statements are its $A=0$
reading. (ii) Proposition~15(b) is the Hellmann--Feynman lift bound; its
worst-case integral $\Sigma(s)=\sigma_0+as$ is replaced by the two-sided
trajectory bracket of Theorem~\ref{thm:bracket} on every eigenvalue of
$PKP$. (iii) The efficient-supply closure of the earlier draft must be dropped: along training the
Rayleigh quotient of the kernel on the rough residual is not a fixed fraction
of $\sigma_R^2$ but decays, $c(t)\approx\lambda_{\mathrm{eff}}(t)/\sigma_R(t)^2\to0$,
because the lift raises the top of the rough subspace while the residual moves
to its slowest directions, the front [n=200: $c=0.02$--$0.09$ at $t=0$,
$0.001$--$0.005$ late, with $\sigma_R^2$ lifted $6$--$50\times$]. The two-sided band
built on it is replaced by the cutoff bound, whose only remainder is the
leakage. (iv) The error is not held in any fixed roughness seminorm: the kernel
lifts the Laplacian-rough directions and fits them, the whole residual becomes
Laplacian-rough while decaying, and what is held is the energy on the modes
still unfit in the \emph{current} basis. Roughness must be measured against
the lifted kernel, which is what the moving frame does. (v) The a~priori
cutoff of Definition~\ref{def:cutoff} becomes the transported cutoff read from
the lifted spectrum; Proposition~\ref{prop:cutoff-floor} remains a valid
a~priori floor. (vi) Proposition~20 is the radius theorem at a single
expansion point and the chain is its continuation. (vii) The reduced model of
ignition is subsumed by the richness expansion and the crossing condition and
can be cut. (viii) The moving frame speaks of rank in the current kernel
basis; the statement that this is roughness in the input-space sense of
Part~I rests on the measured law that the Laplacian roughness of the $k$-th
eigenvector is the same function of $k$ for $K_0$ and $K_T$ [n=200], which has
no proof and must be re-established at the sample sizes of interest, with a
held-out sample, before the cutoff bound is read as a roughness statement.
(ix) The a~priori floor of Theorem~\ref{thm:rough-apriori} has not been
evaluated against $\sigma_R(t)$ on any run; the analogous $H_2$ bound on the
propagator is loose by $e^{13}$ on the reference run, and if the floor is loose
by a comparable factor its content at the times of interest is nil, in which
case the trajectory bracket of Theorem~\ref{thm:bracket} on $PKP$ replaces it.

\subsection{Why the late stage is not predictable from the initialisation}
\label{ssec:chain-impossible}

\begin{proposition}[Determinism; no chaos]
\label{prop:nochaos}
For analytic $N_\theta$ the gradient-flow trajectory is real-analytic in $t$, so
the history on any interval determines $E(T)$. The loss is a strict Lyapunov
function, so there are no recurrent sets and no sustained exponential
separation of nearby trajectories; what remains is finite-time stretching
along transiently negative-curvature directions and algebraic amplification
through the lift. [n=200: over the full run a $10^{-7}$ relative perturbation of $\theta_0$ grows
by a factor $6$ (random direction) and $21$ (tail-feature direction), while a
top-feature perturbation is contracted $16\times$ during the fit of the top mode
and regrows to $5\times$ through the lift; the induced residual perturbation
grows by two to three orders of magnitude and ends at $10^{-6}$ relative. The
growth is algebraic throughout, roughly $t^{0.2\text{--}0.4}$ after $t\approx10$,
with no exponential phase in $3.5$ decades of time.]
\end{proposition}

So the obstruction is conditioning, and it has three parts, each a theorem.

\begin{theorem}[Hidden state]
\label{thm:hidden}
Let $A$ and $\Delta$ be the input activations and the outgoing backprop
signals of a hidden-hidden block on the data at $\theta_0$. For every
$N=UV^{\!\top}$ with $U^{\!\top}\Delta=0$ and $V^{\!\top}A=0$ (such $N$ exist whenever the
widths exceed the ranks, e.g.\ $H>n\,n_y$, and have any norm), the
initialisation $\theta_0+N$ has the same unit states, output, residual, loss,
gradient and kernel on the data, and differs only in the propagation part of
Proposition~\ref{prop:chain-identities}(d). Hence $E(T)$ is not a function of
$(K_0,E_0,Y)$, nor of any finite number of their time derivatives at $0$ when
the widths allow the corresponding null spaces. [Verified: with
$\|N\|_F=\|W\|_F$ the states agree to $5\cdot10^{-15}$ and the kernel increments
after one step differ by $1\%$.]
\end{theorem}

\begin{proposition}[Visibility]
\label{prop:visibility}
Let $B$ be the modes of $K(t_0)$ carrying a fraction $\phi$ of $\|E(T)\|^2$. Their
share of the loss-decrease rate at $t_0$ is at most
$\max_B\lambda_k(t_0)/\lambda_{\mathrm{eff}}(t_0)$, and in the fixed-basis
solution their components are constant on $[0,t_0]$ to relative order
$\lambda_k(t_0)t_0$. [n=200: median $1.4\cdot10^{-6}$ at $t_0=0.01T$ and
$10^{-2}$ at $t_0=0.1T$; share of the rate $8\cdot10^{-7}$ and $7\cdot10^{-4}$.]
\end{proposition}

\begin{proposition}[Amplification]
\label{prop:amplification}
The same modes are later lifted by the gain of
Theorem~\ref{thm:bracket} (measured $10^2$--$10^4$), coherently across blocks,
quadratically in the accumulated drive, and $E_k(T)\sim\exp(-\!\int\lambda_k)$ is
exponentially sensitive to the lifted eigenvalue. In the fixed-basis model the
sensitivity of $\log|E_k(T)|$ to a mode's growth exponent is
$2\Lambda_k(T)\log(T/t_0)$ with $\Lambda_k(T)\approx1$ for the modes that
matter.
\end{proposition}

\begin{corollary}[The conditioning statement]
\label{cor:conditioning}
To predict $E(T)$ to relative precision $\varepsilon$ from the history up to
$t_0$ the tail components must be known to relative precision of order
$\varepsilon\,\lambda_k(t_0)t_0$ divided by the amplification, i.e.\ about
$10^{-7}$ at $t_0=0.01T$ for $\varepsilon=0.1$ [n=200]: five orders of magnitude
below the initialisation-to-initialisation variation of the early history
($2\cdot10^{-2}$) and below single precision. The late state is a smooth,
deterministic function of the early state whose relevant derivatives are
enormous exactly in the directions that are unobservable early.
\end{corollary}

The three parts have a common root: the drive is the residual, the residual
lives in the slow modes of the current kernel, and the slow modes are both the
ones that have not moved yet (invisible) and the ones feature learning
targets (amplified). Feature learning is the mechanism by which the network
amplifies precisely the information that no early observable contains, and
that information is seeded by the initialisation through
Theorem~\ref{thm:hidden}. This is why every early predictor we tried (frozen
kernel, gain-law extrapolation, residual front, the loss itself) orders
designs within $1.1\times$ at chance and within $1.3\times$ at $65$--$76\%$, why
seed pairs that agree to $2\%$ for the first tenth of training end $15\%$
apart, and why the only object that reproduces the late error is the
trajectory carried forward [n=200].

\subsection{What a statistical description can promise}
\label{ssec:chain-statistical}

The impossibility above is a statement about one network. A statistical
description does not predict one network; it predicts the distribution over
initialisations, and by Theorem~\ref{thm:hidden} that is the correct object,
because the quantity that defeats prediction is the draw of the hidden state.
Three consequences follow, two of them proved and one to be measured.

\emph{The irreducible uncertainty of any statistical prediction for a single
run is the spread of the late error over initialisations}, $15\%$ median here
[n=200]. A design comparison finer than that spread is not well posed for one
run of each design, whatever the method; it is well posed for the expected
error, which is a statistic.

\emph{The coherent chain is the statistical approximation}: replace the
network's own units by draws from the unit distribution and the propagation
part by its average, and Propositions~\ref{prop:chain-identities}(b,c) and
Theorem~\ref{thm:spike}(a) close on the data. On the one case run to
convergence [H=96, depth 2, n=200] it reproduced the structure, the gain law
and the ordering of designs, with a $2\times$ bias in the late error and a
$4$--$5\times$ excess in the tail gain traced to the averaged propagation term and
to the transfer of the ensemble's lift; at depth~3 the bias was larger
($5\times$ in the tail gain). It was not run at larger width. Whether the
bias and the seed spread decrease with width, as the propagation term being
a finite-width fluctuation suggests, or persist with depth, as the depth-3
result warns, is the single most consequential unmeasured quantity in this
work, and it decides the question in the next paragraph.

\emph{The diagnostic that is well posed is the expected late-stage ability of an
architecture.} If the spread over initialisations shrinks with width and the
coherent chain converges to the exact one in that limit, then ranking
architectures by their expected late error is both meaningful and computable
without the particular initialisation, and the finite-width exponent of the
convergence is the error bar. Everything in this section is set up for that
measurement: the exact chain for the per-initialisation truth, the coherent
chain for the expectation, and the seed spread for the floor. The ranker
built from the identities (lazy error, $H_2$, top eigenvalue, output scale)
reached a $1.3\times$ resolution out of sample [n=200], which is the expected
accuracy of a one-Jacobian statistic and should be read as a lower bound on
what the expected-error ranking can do.

\subsection{Scope and what to measure}
\label{ssec:chain-scope}

Every number in this section comes from $n=200$ training points on the
Lorenz-63 one-step map, at which the training set is nearly fit within the
budget, the residual-adapted late phase is short, the depth-$6$--$8$ ``floor'' is a
sample-size artefact, and the negative lazy-transfer sign in the validation
split is overfitting of a $200$-point sample. The theorems carry $n$ as a symbol
and are unaffected; the verifications, bulk ratios, exponents, signs, transfer
fractions and the ranker's fits are all $n=200$ statements and are to be redone
at a sample size that leaves a genuine late stage before any is quoted. The
exact chain does not scale to large $n$ (one $N\times N$ eigendecomposition per
stage); the ground-truth runs do, so the redo is training-only, and it should
include, at that $n$: the seed spread of the late error against width and
depth, with depth one as the control in which the hidden-state channel of
Theorem~\ref{thm:hidden} is absent; the gap between the own-initialisation and the other-initialisation
chain against width and depth; the per-stage sign of the increments as the
smallest eigenvalue of their compression to the residual's subspace; the spike bound along
the trajectory; the transported cutoff bracket with its leakage integral; the
rank-to-roughness law with a held-out sample; the a~priori floor against the
measured $\sigma_R(t)$; the validation split; and the Euler-versus-direct comparison
at several $\Delta$ from a shared initialisation.

\paragraph{Hyperparameter transfer (negative result, n=200).} A sweep over widths
$48$--$384$ and output scales $0.25$--$2$ at fixed budget $3000/s^2$ did not
bracket a training-error optimum (the richest scale was best at every width and
depth) and the feature-to-lazy ratio of the descent rate did not collapse across
widths under either the standard or the $\mu$P-like labelling; at fixed scale it
grew with width. On training error at fixed budget, richness has no cost here,
so the transfer question needs a criterion with an optimum (test error or a
stability constraint) and is left open.
